\documentclass[11pt,letterpaper]{article}
\usepackage[margin=1in]{geometry}
\usepackage{fontspec}
\setmainfont{Latin Modern Roman}
\newfontfamily\chinesefont{Noto Serif CJK SC}
\usepackage{amsmath,amssymb,amsthm,mathtools}
\usepackage{booktabs,array,enumitem}
\usepackage[hidelinks,unicode]{hyperref}
\usepackage{microtype}
\setlength{\emergencystretch}{2em}
\setlist{itemsep=2pt,topsep=5pt}
\newtheorem{theorem}{Theorem}[section]
\newtheorem{proposition}[theorem]{Proposition}
\newtheorem{corollary}[theorem]{Corollary}
\newtheorem{lemma}[theorem]{Lemma}
\theoremstyle{remark}
\newtheorem{remark}[theorem]{Remark}
\DeclareMathOperator{\Aut}{Aut}
\newcommand{\mi}{\operatorname{mis}}
\newcommand{\ind}{\mathcal I}
\newcommand{\V}{V}
\newcommand{\N}{N}
\newcommand{\abs}[1]{\lvert #1\rvert}
\title{Independent set compression and Hall extremizers\\
in the sixth Mycielski graph}
\author{}
\date{5 October 2026}
\begin{document}
\maketitle
\vspace{-1.5em}
\begin{abstract}
We give a computer-assisted determination of the Hall ratio of the
47-vertex Mycielski graph, with the convention $M_2=K_2$ and
$M_{r+1}=\mu(M_r)$. Its Hall ratio is $10/3$. Exactly 1,990 induced vertex
subsets attain this value, forming 199 orbits under the full ambient
automorphism group, each of size ten. Every maximizer has twenty vertices,
independence number six, and contains the final apex. A general
neighborhood-union identity generates maximal independent sets of
$\mu(G)$ from independent sets of $G$. For $M_6$, this yields 857 constraints
from 7,407 independent sets of $M_5$. Exact rational certificates establish
the upper bound and equality reductions; two different exact algorithms
independently enumerate all extremizers. The accompanying programs and
certificates avoid enumerating all $2^{47}$ induced subsets.
\end{abstract}

\section{The finite result and its scope}
For a nonempty finite graph $H$, its Hall ratio is
\[
 \rho(H)=\max_{\varnothing\ne S\subseteq V(H)}
 \frac{\abs S}{\alpha(H[S])}.
\]
Here $\alpha$ denotes the independence number, and $H[S]$ is the induced
subgraph. We abbreviate $\alpha(H[S])$ to $\alpha(S)$ when the ambient graph
is clear. We use the ordinary Mycielski construction: $\mu(G)$ has original
vertices $V(G)$, an independent clone layer $V(G)'$, and an apex $z$.
Original vertices retain their edges; $u$ is adjacent to $v'$ precisely
when $uv\in E(G)$; and $z$ is adjacent to every clone. Set $M_2=K_2$ and
$M_{r+1}=\mu(M_r)$, so the successive orders are $2,5,11,23,47$.

\begin{theorem}[Computer-assisted finite classification]\label{thm:main}
The graph $M_6$ satisfies $\rho(M_6)=10/3$. Every maximizing subset has
order $20$, independence number $6$, and contains the final apex.
There are exactly $1,990$ labelled maximizing subsets, in $199$ orbits
under $\Aut(M_6)\cong D_5$. Every orbit has size $10$. The final
original/clone/apex layer counts are as follows.
\begin{center}
\begin{tabular}{@{}ccc rr@{}}
\toprule
Originals & Clones & Apex & Orbits & Labelled subsets\\
\midrule
13 & 6 & 1 & 109 & 1,090\\
14 & 5 & 1 & 83 & 830\\
15 & 4 & 1 & 7 & 70\\
\midrule
\multicolumn{3}{l}{Total} & 199 & 1,990\\
\bottomrule
\end{tabular}
\end{center}
\end{theorem}

The orbits in this theorem concern embedded subsets under automorphisms
of the ambient graph. They are not asserted to be distinct abstract
isomorphism types of the induced graphs.

Cropper, Gy\'arf\'as and Lehel~\cite{CGL} give the values
$\rho(M_2)=2$, $\rho(M_3)=5/2$, $\rho(M_4)=8/3$ and $\rho(M_5)=3$, and prove
that the sequence of Hall ratios is unbounded. The cardinality-profile
formulation used below predates this calculation; see Barnett~\cite{Barnett}.
Steiner~\cite{Steiner} records that boundedness of
$\chi_f(M_r)/\rho(M_r)$ remains open. The present finite calculation does
not resolve that asymptotic question. This note records reproducible
results and complete proofs of the structural reductions; it makes no
claim that publication priority has been established.

\section{Neighborhood compression for a general base graph}
Throughout this section, $G$ is a finite simple graph with $n\ge1$ vertices.
Let $\ind(G)$ contain all independent sets, including the empty set, and
write $i(G)=\abs{\ind(G)}$ and $\mi(G)$ for the number of maximal independent
sets. For $I\subseteq V(G)$, $N(I)$ is the union of its open neighborhoods.
Define
\[
 \mathcal D(G)=\{N(I):I\in\ind(G)\},\qquad d(G)=\abs{\mathcal D(G)}.
\]

\begin{theorem}[Neighborhood-union identity]\label{thm:compression}
For every nonempty finite simple graph $G$,
\begin{equation}\label{eq:count}
 \mi(\mu(G))=\mi(G)+d(G).
\end{equation}
For $N\in\mathcal D(G)$ put $J_N=\{v:N(v)\subseteq N\}$. The maximal
independent sets avoiding the apex are exactly
\[
 K_N=J_N\cup(V(G)\setminus N)',\qquad N\in\mathcal D(G),
\]
with one distinct set for each $N$. Those containing the apex are exactly
$\{z\}\cup I$, where $I$ is maximal independent in $G$.
\end{theorem}
\begin{proof}
Choose independent $I$ with $N=N(I)$. Then $I\subseteq J_N$ and
$J_N\cap N=\varnothing$. Indeed, if $v\in J_N\cap N$, it has a neighbor
$i\in I$, and $N(v)\subseteq N$ would force $i\in N$, contrary to the
independence of $I$. It follows that $J_N$ is independent and
$N(J_N)=N(I)=N$.

Thus $K_N$ is independent. An excluded original vertex has a neighbor
among its selected clones, by the definition of $J_N$. An excluded clone
has a neighbor in $I\subseteq J_N$. There is a selected clone: $V\setminus N$
contains $I$ when $I$ is nonempty, and equals $V$ otherwise. Hence the apex
also has a neighbor in $K_N$, so $K_N$ is maximal.

Conversely, a maximal independent set avoiding $z$, with original part
$I$, must contain every available clone, namely $(V\setminus N(I))'$.
Maximality then forces its original part to be $J_{N(I)}$. The clone part
uniquely determines $N$. A maximal independent set containing $z$ contains
no clone, and its original part is maximal in $G$. These two bijections
prove~\eqref{eq:count}.
\end{proof}

\begin{corollary}\label{cor:complexity}
All maximal independent sets of $\mu(G)$ can be generated in
$O(n^2i(G))$ elementary Boolean operations, with
$O(n^2+ni(G))$ bits of working and output storage. Their number is at most
$2i(G)$.
\end{corollary}
\begin{proof}
Enumerate independent sets of $G$, collect their distinct neighborhood
unions, and identify maximal sets by $I\cup N(I)=V(G)$. Construct the sets
in Theorem~\ref{thm:compression}. Boolean adjacency arrays give the stated
bounds; bitsets improve the implementation. Finally,
$d(G)\le i(G)$ and $\mi(G)\le i(G)$.
\end{proof}

This is a bound in the number of independent sets of the base graph.
It does not imply polynomial-time Hall optimization in the number of
vertices, or a polynomial bound on the size of a branch certificate.
For the present instance, exact enumeration gives
\begin{equation}\label{eq:counts}
 i(M_5)=7,407,\qquad d(M_5)=778,\qquad \mi(M_5)=79,
 \qquad \mi(M_6)=857.
\end{equation}

\begin{proposition}[Induced-subset recurrence]\label{prop:recurrence}
Let $A,B\subseteq V(G)$ and $\varepsilon\in\{0,1\}$. For the induced subset
consisting of originals $A$, clones $B'$, and the apex when
$\varepsilon=1$, its independence number equals
\begin{equation}\label{eq:recurrence}
 \max\left\{\varepsilon+\alpha(G[A]),\quad
 \max_{\substack{I\subseteq A\\ I\in\ind(G)}}
 \bigl(\abs I+\abs{B\setminus N(I)}\bigr)\right\}.
\end{equation}
\end{proposition}
\begin{proof}
An independent set avoiding the apex chooses independent original part
$I$ and may include every available clone. A set using the apex chooses
no clone and may include a maximum independent set in $G[A]$. If
$\varepsilon=0$, the first term is redundant and does not change the
maximum.
\end{proof}
The same decomposition gives the exact independence-polynomial identity
\[
 Z_{\mu(G)}(t)=tZ_G(t)+\sum_{I\in\ind(G)}t^{\abs I}
 (1+t)^{n-\abs{N(I)}},\qquad
 Z_G(t)=\sum_{I\in\ind(G)}t^{\abs I}.
\]
For example, it yields $i(M_6)=39,473,983$ without enumerating those sets.

\section{An exact certificate for the Hall ratio}
Let $K_1,\ldots,K_q$ be all maximal independent sets of a graph $H$, and
let $x_v\in\{0,1\}$ indicate a subset $S$. Then
\begin{equation}\label{eq:profile}
 \alpha(H[S])\le k
 \quad\Longleftrightarrow\quad
 \sum_{v\in K_j}x_v\le k\quad(1\le j\le q).
\end{equation}
Every independent set extends to a maximal one, which proves the reverse
implication; each intersection $S\cap K_j$ proves the forward one.
Thus maximizing $\sum_vx_v$ under~\eqref{eq:profile} is an exact binary
model for the cardinality profile. For $M_6$ it has 47 variables and
857 independent-set rows.

\subsection{Rational leaves and binary splits}
At a certificate node, fix vertices $O$ to one and $Z$ to zero, leaving
$U=V(H)\setminus(O\cup Z)$ free. Set $b_j=k-\abs{K_j\cap O}$.
A rational-dual leaf supplies $y_j\ge0$ and $z_v\ge0$ satisfying
\begin{equation}\label{eq:coverage}
 \sum_{j:v\in K_j}y_j+z_v\ge1\qquad(v\in U).
\end{equation}
Multiplying the residual rows by $y_j$ and the bounds $x_v\le1$ by $z_v$
proves
\begin{equation}\label{eq:dual}
 \abs S\le B:=\abs O+\sum_j y_jb_j+\sum_{v\in U}z_v.
\end{equation}
If the exact rational $B$ is below an excluded integer cardinality $T$,
that node contains no subset of order at least $T$. An internal node
splits on a free binary variable and checks both assignments. The
certificate therefore covers every subset. Checking the displayed
inequalities uses exact integer and rational arithmetic; numerical
optimization is used only to propose weights or branches.

\subsection{Upper bound and witness}
For $k\in\{1,3,4,\ldots,14\}$, the supplied certificate excludes
\[
 T_k=\left\lfloor\frac{10k}{3}\right\rfloor+1.
\]
The case $k=5$ uses 130 splits and 131 rational leaves. Each of the other
twelve cases uses one leaf, for 273 nodes in total. The missing case
$k=2$ is elementary: every six vertices of a triangle-free graph contain
an independent triple. Given a vertex, three neighbors are independent;
otherwise it has three nonneighbors, among which a nonedge together
with the vertex gives an independent triple. Thus $\alpha\le2$ implies
$\abs S\le5$. For $k\ge15$, use $\abs S/k\le47/15<10/3$.
The graph is checked to be triangle-free, so these cases prove
$\rho(M_6)\le10/3$.

For the reverse bound, label originals $0,\ldots,n-1$, their clones
$n,\ldots,2n-1$, and the apex $2n$ at each step. The subset
\begin{align*}
 S=\{&0,1,2,3,4,5,6,7,11,12,15,19,21,22,\\
     &26,32,33,41,45,46\}
\end{align*}
has twenty vertices and independence number six. Its maximum
intersection with all 857 maximal independent sets is exactly six;
$\{0,2,5,7,11,41\}$ is one independent six-set. The independent audit also
checks its independence number by a direct include/exclude recursion on
these twenty vertices. Hence $\rho(M_6)=10/3$.

\section{Classifying all equality subsets}
Equality forces $\alpha(S)$ to be divisible by three. The only possible
pairs are
\[
 (10,3),\quad(20,6),\quad(30,9),\quad(40,12).
\]
Additional exact certificates
exclude the first, third and fourth pairs using 19, 1 and 1 nodes,
respectively. One further leaf excludes a twenty-set with independence
number at most six when the final apex is absent. These 22 nodes have
nine binary splits and thirteen rational leaves.

Consequently every maximizing subset has twenty vertices, independence
number six, and contains the apex. Write it as $A\cup B'\cup\{z\}$.
The apex implies $\alpha(M_5[A])\le5$, while $\abs B\le6$. Exact base
independence values give
\begin{equation}\label{eq:baseprofile}
 \max_{\alpha(M_5[A])\le4}\abs A=12,\qquad
 \max_{\alpha(M_5[A])\le5}\abs A=15.
\end{equation}
These bounds are reverified by both complete enumerations. Since
$\abs A+\abs B=19$, they imply
\[
 \abs A\in\{13,14,15\},\quad \alpha(M_5[A])=5,
 \quad \abs B=19-\abs A.
\]
By Proposition~\ref{prop:recurrence}, such a pair is a maximizer exactly
when
\begin{equation}\label{eq:family}
 \abs{B\setminus N(I)}\le6-\abs I
 \qquad\text{for every independent }I\subseteq A.
\end{equation}
This is a complete neighborhood description of the three structural
families, not just a condition satisfied by the displayed witness.

\subsection{The full ambient automorphism group}
The graph $M_3$ is $C_5$, in cyclic order $(0,1,2,4,3)$. At each later
stage the apex is the unique maximum-degree vertex. Indeed, for a base
of order $n$ and maximum degree $(n-1)/2$, the new original degrees are
at most $n-1$, clone degrees at most $(n+1)/2$, and apex degree $n$.
The property starts with $C_5$ and inducts.

Every automorphism therefore fixes the apex and preserves both layers.
Its restriction to originals is an automorphism of the base. The bases
are open-neighborhood twin-free, beginning with $C_5$, and this forces
the clone permutation to agree with the original permutation. Conversely,
every base automorphism extends. Thus $\Aut(M_6)\cong D_5$ has order ten.
The independent checker obtains the ten automorphisms of the initial
cycle by testing its 120 permutations, then lifts them and verifies
adjacency, closure, and the unique-apex and twin-free hypotheses.

\subsection{Two complete exact enumerations}
The primary enumeration computes $\alpha(M_5[A])$ for all $2^{23}$ base
subsets by vertex-deletion dynamic programming. It keeps only canonical
$A$ of sizes 13--15 with independence number five. There are 4,370 such
$D_5$ representatives. For each one, choose an independent five-set
$I_0\subseteq A$ minimizing $\abs{N(I_0)}$. Equation~\eqref{eq:family}
requires at most one clone index outside $N(I_0)$. Thus, for
$b=19-\abs A$, enumerate only
\begin{enumerate}
\item $b$ indices from $N(I_0)$;
\item $b-1$ indices from $N(I_0)$ and one from its complement.
\end{enumerate}
Check the remaining inequalities~\eqref{eq:family}, then canonicalize the
full subset under $D_5$. The choice of $I_0$ changes efficiency, not
coverage: every admissible $B$ appears in exactly one of the two cases.

This run tests 50,504,405 filtered clone candidates. It accepts 200 pairs
at canonical $A$ representatives, which become 199 full-subset orbits.
The final deduplication is needed because a stabilizer of $A$ can act
nontrivially on $B$.

A separate implementation constructs all independent sets of $M_5$ and
computes $\alpha(A)$ by a maximum subset-zeta transform, rather than the
primary recurrence. It enumerates clone sets through recursive capacity
propagation of~\eqref{eq:family}, rather than the anchor-neighborhood
filter. It explores 184,588 recursive nodes and independently produces
exactly the same sorted 199 canonical masks. Its completeness proof
checks every pruning rule: saturation, insufficient remaining vertices,
and the residual cardinality bound of a capacity constraint.

Both methods cover every original-set orbit and every clone choice
permitted by~\eqref{eq:family}. All 1,990 labelled subsets in the output
orbits are checked against all 857 maximal independent sets. Every
stabilizer is trivial. The orbit counts in Theorem~\ref{thm:main} follow.

\section{Reproducibility and verification boundary}
The core verification commands are below. Run the first in
\path{mycielski_pilot} and the remaining commands in
\path{mycielski_strengthening}. The accompanying manifests identify all
source files, certificates, orbit data, and independent audit records.
\begin{quote}\begin{minipage}{\linewidth}\small\ttfamily
python check\_exact\_certificate.py\par
 g++ -O3 -std=c++17 enumerate\_extremizers.cpp -o enumerate\_extremizers\par
 ./enumerate\_extremizers\par
 python check\_strengthening.py
\end{minipage}\end{quote}
The independent audit has its own enumeration and checker. Python checking requires version 3.10 or later and only the
standard library; regenerating proposed LP weights additionally uses
SciPy. The primary complete enumeration has a safety time limit: a limit
exit is explicitly incomplete and must not be treated as a proof.

The general structural module is tested on all 75 simple nonempty base
graphs of order at most four, including 33,864 induced-subset recurrence
checks. These finite tests supplement the general proofs. All proof
certificates and enumeration code are preserved with checksums.
No floating-point solver status is a proof premise. No proof-assistant
formalization is claimed, and no conclusion about the asymptotic ratio
$\chi_f(M_r)/\rho(M_r)$ is inferred.

\section*{\chinesefont 中文摘要}
\begingroup\chinesefont
\XeTeXlinebreaklocale "zh"
\XeTeXlinebreakskip=0pt plus 1pt

本文给出第六个 Mycielski 图的 Hall 比的计算机辅助精确证明。采用
$M_2=K_2$、$M_{r+1}=\mu(M_r)$ 的编号约定，$M_6$ 有 $47$ 个顶点，且
$\rho(M_6)=10/3$。达到该比值的诱导顶点集恰有 $1\,990$ 个；在整个图的
自同构群 $D_5$ 作用下，它们分成 $199$ 个轨道，每个轨道含 $10$ 个顶点集。
每个极值集都有 $20$ 个顶点、独立数 $6$，并包含最后加入的顶点。按原顶点、
复制顶点和最后顶点的数量划分，三类轨道数分别为 $109$、$83$ 和 $7$。
一般的邻域并集恒等式把 $\mu(G)$ 的极大独立集生成问题压缩到基图 $G$ 的
独立集枚举。对于本例，$M_5$ 的 $7\,407$ 个独立集生成 $M_6$ 的 $857$ 个
极大独立集约束。精确有理数证书证明上界及等号情形的必要限制，两种不同的
整数算法独立完成全部极值集枚举，避免遍历 $2^{47}$ 个诱导顶点集。本文只陈述
可复核的有限结论，不声称已确立发表优先权，也不解决该图族的渐近问题。

\endgroup

\begin{thebibliography}{9}
\bibitem{CGL}
M.~Cropper, A.~Gy\'arf\'as and J.~Lehel,
\emph{Hall ratio of the Mycielski graphs},
Discrete Mathematics \textbf{306} (2006), 1988--1990.
\href{https://doi.org/10.1016/j.disc.2005.09.020}{doi:10.1016/j.disc.2005.09.020}.

\bibitem{Barnett}
J.~B.~Barnett,
\emph{The Fractional Chromatic Number and the Hall Ratio},
Ph.D. dissertation, Auburn University, 2016.
\href{https://etd.auburn.edu/handle/10415/5316}{Auburn dissertation repository},
Definition~2.1.1 and Chapter~4.

\bibitem{Steiner}
R.~Steiner,
\emph{Fractional Chromatic Number Vs. Hall Ratio},
Combinatorica \textbf{45} (2025), article~37.
\href{https://doi.org/10.1007/s00493-025-00164-0}{doi:10.1007/s00493-025-00164-0}.
\end{thebibliography}
\end{document}
